\section{HDI and ETI: Credible Interval Variants}\label{app:hdi}

Given a posterior density $p(\theta \mid \mathbf{x})$, where $\theta$ is the parameter
of interest and $\mathbf{x}$ the observed data, a $(1-\alpha)$ credible interval
$[L, U]$ satisfies
\begin{equation}\label{eq:credible_interval}
  \int_L^U p(\theta \mid \mathbf{x})\,d\theta = 1 - \alpha.
\end{equation}
Infinitely many intervals satisfy this constraint. This paper uses the Highest Density
Interval throughout; the Equal-Tailed Interval is described below for comparison.

\subsection*{Highest Density Interval (HDI)}

The HDI is the shortest interval satisfying Equation~\ref{eq:credible_interval}:
\begin{equation}\label{eq:hdi}
  [L^*, U^*] = \underset{[L,U]:\;\int_L^U p(\theta|\mathbf{x})\,d\theta\,=\,1-\alpha}{\arg\min}\;(U - L).
\end{equation}
An equivalent characterisation for unimodal posteriors: $[L^*, U^*]$ is the unique
interval for which $p(L^* \mid \mathbf{x}) = p(U^* \mid \mathbf{x})$, so every point
inside has posterior density at least as high as every point outside (hence ``highest
density'').\footnote{For multimodal posteriors, multiple intervals can satisfy the
equal-density boundary condition, and the HDI may be a union of disjoint intervals;
the shortest-interval definition in Equation~\ref{eq:hdi} selects the one with minimum
total width. All posteriors in this paper are unimodal.}
For the Binomial model used in this work, under a flat prior $\mathrm{Beta}(1,1)$, the
posterior after $s$ successes and $f$ failures is $\mathrm{Beta}(s+1,\,f+1)$, and the
HDI is found by minimising the interval width over the lower tail
probability.\footnote{Specifically, minimise
  $F^{-1}(1-\alpha+\ell)-F^{-1}(\ell)$ over $\ell\geq 0$; the minimiser $\ell^*$
  gives $L^*=F^{-1}(\ell^*)$ and $U^*=F^{-1}(1-\alpha+\ell^*)$
  (\texttt{HDIofICDF} in the accompanying code).}

\subsection*{Equal-Tailed Interval (ETI)}

The ETI places $\alpha/2$ probability mass in each tail:
\begin{equation}\label{eq:eti}
  L_{\rm ETI} = F^{-1}\!\left(\tfrac{\alpha}{2}\right), \qquad
  U_{\rm ETI} = F^{-1}\!\left(1 - \tfrac{\alpha}{2}\right),
\end{equation}
where $F^{-1}$ is the posterior quantile function. The ETI is simple to compute and
uniquely defined, but it is not generally the shortest interval.
For symmetric, unimodal posteriors the HDI and ETI coincide; they diverge for skewed
distributions (for instance, a Beta posterior near $\theta = 0$ or $\theta = 1$)
where the HDI is noticeably shorter and therefore preferred.


